\section{Conclusion}
\label{sec:conclusion}

A CVE entry records that a vulnerability exists; it does not let anyone
observe it. We presented \cvexthree, a standard format for executable CVEs
that packages the vulnerable target, a benign workload, an attack, a
recorder, and a machine-checkable success criterion into a self-contained,
container-based record. Inspired by Capture-the-Flag competitions and by the
verifiers of \genie, \cvexthree defines success as the capture of a secret
that only a working exploit can reach, and it runs every record twice, so
that each run yields a labeled pair of benign and attack captures at the
network and system-call level.

We release a corpus of \NumRecords records, \NumFunctional of which are
fully functional and ship with verified reference datasets. We hope that
\cvex will make reproductions durable and shareable artifacts, and that the
corpus will be useful to researchers who build detection models, assess
patches, and evaluate automated security tools. Functionalizing the
remaining records, and extending the format to virtual machines for
vulnerabilities that require OS-level control, are the next steps.
